\subsection{Benefits}

GenAI represents a significant step forward in the development of modern digital technologies. Its applications cover a wide range of fields, from business and the creative industry to education and medicine. By leveraging the capabilities of automatic content generation and data analysis, GenAI contributes to process optimization, increased efficiency, and drives innovation in various areas of human activity.

\subsubsection{Healthcare and Life Sciences [B1]}
\paragraph{Diagnostics, Decision-Making, and Patient Care}
GenAI enhances diagnostic accuracy and clinical decision-making through multiple mechanisms. LLM-based models such as GPT-4 integrate extensive medical literature to recommend relevant examinations and provide evidence-based diagnostic insights \parencite{S01}, while also improving medical imaging for disease detection and monitoring \parencite{R09}. NLP solutions automatically extract data from medical documents, identify patient conditions and disease severity, and strengthen interdepartmental communication and care coordination \parencite{S10}. GenAI models further mitigate diagnostic bias through multimodal data integration and advanced machine learning techniques \parencite{S16}, with applications in triage and diagnosis support \parencite{S02}, retinal healthcare \parencite{S03}, dental telemedicine, and chest radiograph interpretation \parencite{S17}. The technology also delivers high-quality semantic health information \parencite{S11} and supports the development of personalized treatment plans through the analysis of large patient datasets \parencite{S16}.

\paragraph{Drug Discovery Acceleration and Precision Medicine}
GenAI accelerates drug discovery by predicting drug properties and supporting genomic research \parencite{S01}. It enhances the design and synthesis of novel compounds, expands and optimizes compound libraries, and enables the creation of molecules with targeted therapeutic properties \parencite{S16}. By reducing the human, material, and financial resources typically required in traditional drug development, GenAI advances precision medicine through the rapid identification and optimization of candidate drugs tailored to specific patient populations and disease characteristics.

\paragraph{Clinical Documentation and Workflow}
GenAI demonstrates enhanced efficiency in medical documentation through literature synthesis and data organization, handling diverse data types including textual patient records, diagnostic reports, research papers, medical imaging data such as MRIs and CT scans, voice recordings, and biomarkers \parencite{S01}. ChatGPT offers substantial potential for automating administrative tasks \parencite{S03} and supporting clinical documentation processes \parencite{S02}, while NLP-based technologies improve interdepartmental communication, coordinate patient care, and ensure appropriate follow-up care \parencite{S10}.

According to \citeauthor{S05}, research reports significant efficiency gains, with AI systems significantly reducing documentation time, particularly in ambient intelligence systems and complex clinical cases \parencite*{S05}. These time savings directly impact clinician workload and patient care availability, offering a promising means to reduce healthcare professional burnout. Six of nine studies reviewed found that AI-generated documentation met or exceeded traditional standards. Healthcare professionals also reported improved usability and reduced cognitive load, supporting broader adoption of AI-assisted documentation. Additionally, NLP-based systems function as virtual assistants for health professionals, streamlining both clinical and administrative workflows \parencite{S10}.

\paragraph{Strengthening Mental Health}
GenAI demonstrates strong potential to enhance mental health by providing accessible emotional and psychological support. Chatbots and conversational interfaces offer patients guidance and reassurance, delivering benefits across educational, healthcare, and broader social contexts \parencite{R09}. GenAI facilitates preliminary patient consultations and psychological assistance, helping patients manage the psychological stresses associated with illness \parencite{S01}. Research highlights its effectiveness in improving holistic understanding, reducing workload for mental health professionals, mitigating loneliness, and reducing the emotional burden on patients \parencite{S17}, collectively contributing to better mental health outcomes and quality of life through scalable, accessible support systems.

\paragraph{Access to Healthcare Knowledge and Simplified Communication}
GenAI improves access to healthcare knowledge and facilitates the communication of complex medical information, enabling patients to obtain reliable insights more quickly and make informed decisions \parencite{R09}. It simplifies medical terminology and statistical data, providing patients with foundational knowledge before consultations and enhancing their understanding of medical results \parencite{S01}. GenAI also strengthens doctor–patient interactions through preliminary consultation tools and clearer explanations, while applications such as ChatGPT show strong potential for patient education \parencite{S01,S02,S03}.

NLP-based technologies act as virtual assistants, offering patients information and support with tasks such as planning, follow-up, and scheduling \parencite{S10}. They provide high-quality, semantically rich health information by simplifying medical texts, conveying disease information effectively, and addressing low-risk health queries \parencite{S11}. GenAI further enhances accessibility by translating medical reports into plain language, generating personalized health guidance, and supporting lifestyle interventions \parencite{S14,S17}. In dental telemedicine, its multilingual capabilities improve scalability and facilitate effective consultations \parencite{S17}, helping democratize access to healthcare knowledge across diverse populations and literacy levels.

\subsubsection{Education and Learning [B2]}
\paragraph{Personalized Learning}
GenAI supports personalized learning through digital teaching assistants and the creation of supplemental materials such as teaching cases and recap questions \parencite{R04}. In medical education, it enables advanced training with real-time simulations \parencite{S01} and acts as a virtual assistant that generates educational content and personalized study plans \parencite{S10}. ChatGPT further adapts content delivery to individual needs, fostering active engagement, self-paced learning, and deeper understanding of the subject \parencite{S12}.

GenAI also promotes equitable education by providing flexible, efficient and cost-effective learning opportunities. It offers instant feedback and explanations that improve self-directed learning and curiosity \parencite{S12}, while facilitating rapid information access and innovative teaching approaches \parencite{S14}. AI-driven tools, such as chatbots, enhance these experiences by supporting creativity, automation, personalization, collaboration, multimodal content creation, and accessibility \parencite{S15,S18}. Collectively, GenAI advances educational systems that adapt to individual learner profiles, preferences, and learning trajectories in diverse contexts.

\paragraph{Digital Learning Resources, Reduced Educator Administrative Tasks}
GenAI streamlines the creation of digital learning resources while significantly reducing the administrative workload of educators. It enables the development of diverse and engaging materials that accommodate different learning styles and enhance instructional quality \parencite{S12}. GenAI reduces the time spent on routine tasks by automating tasks such as generating multiple-choice questions, planning lessons, and supporting technology-based teaching \parencite{S12}. It also assists in creating new educational content and exam questions, with automatic scoring and grading capabilities \parencite{S14}, allowing educators to focus on higher-value instructional activities and student engagement.

\paragraph{Assistance in Language Translation and Accessibility}
GenAI enhances accessibility and inclusivity in education by supporting language translation and adapting content for diverse learner populations. It streamlines classroom tasks such as lesson planning and technology-based instruction while providing instant answers that promote self-directed learning and exploration \parencite{S12}. GenAI also assists in writing assignments, developing research papers, and generating educational materials and exam questions \parencite{S14}, ensuring that learning resources are accessible to students with different linguistic and accessibility needs, including those with disabilities.

\subsubsection{Research, Innovation, and Design [B3]}
\paragraph{Supports Design Knowledge and Research Across Disciplines}
GenAI demonstrates transformative potential in improving productivity, decision-making, and economic value across business sectors and research disciplines \parencite{R01}. It improves efficiency \parencite{R03} and supports process discovery by generating process descriptions that help organizations identify and analyze different workflow stages \parencite{R04}. The capacity of GenAI to model complex, non-linear business processes enables its use in implementation, simulation, and predictive monitoring, while fostering innovation through new business ideas, products, services, and models \parencite{R04}.

The technology reshapes organizational knowledge management by automating knowledge discovery from large volumes of unstructured, distributed data. It enhances knowledge sharing through automatic generation and dissemination of multilingual content, such as Wikis and FAQs, and delivers personalized insights to employees \parencite{R04}. In design science research, GenAI supports the construction of novel IT artifacts by extracting design knowledge in the form of requirements, principles, and features, from interdisciplinary sources, making it collectively available to researchers and practitioners \parencite{R04}. Integrated into design thinking and related methodologies, GenAI augments human creativity in idea generation, user needs elicitation, prototyping, evaluation, and automation \parencite{R04}. Furthermore, it enables the algorithmic identification of knowledge gaps and inconsistencies, promotes new dialogic and methodological approaches, and supports the formulation of innovative research questions across disciplines \parencite{R06,S02}.

\paragraph{Tool-Supported Idea Generation}
According to \citeauthor{R04}, GenAI enhances idea generation across organizational functions by combining human and computational creativity \parencite*{R04}. In business process management, it supports innovative process redesign and automation, driving the development of next-generation process guidance systems. It enables automated knowledge discovery, improves knowledge sharing through content generation, and maintains enterprise models at multiple abstraction levels, while supporting digital twin applications for enterprise asset management.

GenAI also delivers high-quality natural language interfaces that enhance usability and accessibility, producing optimized content for social media, emails, and reports \parencite{R04}. It improves collaboration through intelligent agents, automates personalized marketing, and strengthen recommender systems through advanced personalization \parencite{R04}. In design thinking and innovation contexts, GenAI supports user needs elicitation, prototyping, evaluation, and design automation \parencite{R04}, showcasing strong potential for human–AI collaboration that amplifies creative problem-solving \parencite{R05,R06}. In architecture, engineering, and construction, it facilitates concept visualization and generation of alternative design solutions, supports data-driven decision-making and provides instant training and feedback \parencite{S15}. GenAI has demonstrated exceptional creative performance, including passing university-level exams, achieved through reinforcement learning from human feedback \parencite{S18,R07}.

\subsubsection{Software Engineering and Technical Productivity [B4]}
\paragraph{Analysis, Coding, Testing and Translations (Automatically)}
GenAI provides enhanced support for content analysis and code generation \parencite{R05}, enabling developers to automate routine coding tasks and improve development efficiency. It generates and optimizes test cases to accelerate testing and meet coverage criteria \parencite{S06}. Additionally, GenAI and large language models enable seamless code translation between programming languages, reducing manual effort and supporting more efficient, interoperable software development workflows \parencite{S06}.

\paragraph{Efficiency in Software Development}
GenAI improves the efficiency of software development by improving productivity, optimizing resource utilization, and reducing operational costs \parencite{R02}. Its integration supports strategic decision-making and fosters human–AI collaboration that augments creative problem-solving within development teams \parencite{R05}. Advanced tools such as Bard, ChatGPT, and Copilot contribute to the design of more accurate and robust software, enabling the production of higher-quality systems in shorter development cycles \parencite{S06}. Furthermore, by incorporating pair programming methodologies derived from Extreme Programming, AI agents can function as collaborative team members, assisting developers throughout the software lifecycle to accelerate time-to-market and enhance code quality and maintainability \parencite{S06}.

\subsubsection{Communication, Accessibility, Services, and Social Impact [B5]}
\paragraph{Communication and Accessibility}
GenAI facilitates the bridging of communication gaps and the delivery of tailored services across diverse audiences, promoting societal inclusion through enhanced engagement and mutual understanding. Intelligent automation enables organizations to provide personalized and adaptive services at scale, resulting in favorable outcomes for more people in society \parencite{R08}. Integrating personalization and automation into complex processes enables GenAI to produce customized content and interactions that support informed decision-making and address individual needs \parencite{R09,S09}.

Moreover, the technology fosters social cohesion by improving cross-cultural and interdisciplinary communication, thereby enhancing societal connectivity \parencite{S09}. As GenAI redefines traditional workflows and user interactions, its integration requires adaptive socio-technical frameworks that reflect evolving modes of human–AI collaboration \parencite{S13}. It offers distinct advantages, including creativity, automation, personalization, multimodal content generation, and improved accessibility, while ensuring responsible adoption requires the use of explainable GenAI approaches \parencite{S18}.

\paragraph{Enhanced User Interaction and Experience}

GenAI facilitates more natural, efficient, and adaptive communication between users and systems, rendering products and services increasingly intuitive and personalized. When integrated with LLMs, smart devices, and the Internet of Things, GenAI functions as an intelligent assistant that enhances individual support, productivity, and overall user experience \parencite{R02}. It further enables the automated generation of personalized marketing content tailored to personality traits, such as introversion or extroversion, demonstrating superior effectiveness compared to uniform communication strategies \parencite{R04}.

Within service marketing and customer relationship management, GenAI supports strategic planning and operational execution by streamlining service delivery and improving customer engagement \parencite{R08}. It facilitates the design of personalized service offerings and the development of targeted marketing strategies for specific customer segments, enabling scalable and adaptive service personalization through intelligent automation \parencite{R08}. In parallel, NLP capabilities enable automated data extraction and analysis, while GenAI-driven conversational platforms simulate human-like interaction, contributing to widespread adoption across diverse domains \parencite{S10}. As these technologies continue to reshape user interactions and organizational workflows, their integration requires flexible socio-technical frameworks that accommodate the evolving patterns of human–AI collaboration \parencite{S13}.

\paragraph{Digital Government Services}

GenAI enhances digital government services through translation and accessibility features that increase service reach and improve citizen engagement. It improves efficiency across the public sector \parencite{R03}, supporting the digital management of non-tangible organizational assets, such as procedures, legal texts, and service documentation, throughout their lifecycles. Comparable advantages extend to the management of physical assets in Industry 4.0 environments \parencite{R04}.

In digital service delivery, GenAI improves the performance of existing services by producing human-like conversations with customers, providing personalized and cost-effective services \parencite{R08}. It serves as a disruptive force across digital services including video streaming, recommendation agents on e-commerce platforms, online financial and banking services, education, legal, and healthcare services \parencite{R08}. Governments can leverage GenAI’s translation and text-to-speech technologies to broaden access to public services, while its content moderation and misinformation detection capabilities contribute to safer and more equitable digital ecosystems \parencite{R08}.

\subsubsection{Data Management, Creative Content, Privacy, and Ethics [B6]}
\paragraph{Summaries and Notes}

GenAI automates information summarization and note generation, thereby improving knowledge management and organizational efficiency. It enhances collaboration within teams by providing intelligent agents that suggest, summarize, and synthesize information based on team context, such as through automated meeting notes \parencite{R04}. The technology creates summaries and notes for various applications, including medical contexts such as surgeries \parencite{S01}, enabling knowledge workers to extract essential insights from complex information while reallocating time to higher value analytical and decision-making tasks.

\paragraph{Content Creation}

GenAI accelerates content creation across multiple media formats and enables novel creative applications. It automates various tasks in marketing and media, including news writing, summarization of web content for mobile platforms, thumbnail generation, and accessibility adaptations such as text-to-speech and Braille-supported content \parencite{R04}. Beyond text, GenAI facilitates multimodal content generation encompassing images, audio, and video \parencite{S18}, while reducing labeling requirements and expanding content creation use cases \parencite{S04,R07}.

However, the same capabilities also enable the production of realistic disinformation, including fake news and propaganda, which are increasingly difficult to detect. Advances in GenAI have lowered the cost of disinformation generation and introduced unprecedented personalization by adapting tone and narrative to specific audiences \parencite{R04}. Moreover, GenAI can replace traditional crowdsourcing through automated annotation and execution of knowledge tasks, underscoring the need for robust ethical and governance frameworks to balance creative innovation with responsible information dissemination.

\paragraph{Synthetic Data, Reducing Bias, Increasing Responsibility}

GenAI facilitates the generation of synthetic data to enhance privacy, mitigate harmful biases, and promote ethical and responsible AI practices. Advances in generative models, particularly Generative Adversarial Networks (GANs), have improved the accuracy and realism of synthetic data while maintaining privacy protection \parencite{S07}. In medical and scientific research, GenAI enables the creation of high-quality datasets that preserve patient confidentiality and support data-driven innovation \parencite{S13}.

Beyond research, synthetic data generation supports organizations in improving operational efficiency, reducing costs, and enhancing service delivery in a secure and ethical manner \parencite{R08}. By enabling model training and validation without exposing sensitive information, GenAI contributes to fairer and more transparent AI systems through the creation of balanced datasets that better represent diverse populations and contexts.

\paragraph{Task Automation and Service Scaling}
GenAI automates routine tasks and enables scalable business processes, thereby enhancing productivity and allowing human workers to focus on higher-value activities. It improves organizational efficiency \parencite{R03} by automating key business process management functions such as process extraction from text, event management, resource allocation, and social media operations \parencite{R04}. GenAI further increases productivity by automating content creation, customer service, and code generation, with the potential to transform entire industries through large-scale process optimization \parencite{R04}.

In service contexts, GenAI enhances customer experience through authentic automation and cost-effective personalized interactions \parencite{R08}. It enables scalable and efficient service delivery while reducing employee workload, improving both productivity and job satisfaction \parencite{R08}. By automating complex organizational processes \parencite{R09} and handling routine tasks such as generating budget proposals \parencite{S01} and data analysis \parencite{S14}, GenAI facilitates greater operational agility and informed decision-making. While offering clear benefits in automation, personalization, collaboration, and accessibility, responsible adoption requires transparency and interpretability through explainable GenAI approaches \parencite{S18}.


